\scanpage{234}
\mspage{184}
% L1–2
$\pi_n(X,x)$ par des classes d'homotopie d'applications ponctuées
% L3
\[
(S^n,x_0)\xrightarrow{\varphi}(X,x)
\]
% L4
j'ai bien, pour une telle application $\varphi$,
% L5
\[
\varphi^*:H^n(X,\unsure{G|C})\to H^n(S^n,\overbrace{\varphi^*\deleted{\unsure{(G)}}G_{\deleted{\unsure{x_0}}}}^{\simeq(G_x)_{S^n}})\simeq G_{\unsure{x|x_0}}
\]
% L6–9
(tenant compte que $S^n$ est 1-connexe, et que sa $n$-cohomologie à val. dans un faisceau constant $M$ est $M$. Sûrement l'application
% L10
\[
H^n(X,G)\unsure{-|\to}G_x
\]
% L11–12
ainsi définie ne dépend que de la classe d'homotopie ponctuée de $\varphi$, d'où
% L13
\[
\pi_n(X,x)\times H^n(X\unsure{,}G)\longrightarrow G_x,
\]
% L14–18
et il faut de plus prouver a) que c'est compatible avec les opérations de $\pi=\pi_1(X,x)$, et b) que c'est additif en $\pi_n(X,x)$, pour avoir l'application \unsure{souhaitée|cherchée}
% L19
\[
\deleted{\unsure{f}}H^n(X,G)\unsure{-|\to}\mathrm{Hom}_\pi(\pi_n(X,x),G_x).
\]
% L20–23
Si j'admets ceci, il n'y a qu'un pas pour admettre aussi que la flèche qui nous \unsure{retient|retiendra} \unsure{vraiment} (dans la suite exacte p. 181)
% archive sheet number: 234
